\chapter{One Bluespec Description, Four Profiles}
\label{ch:bsv}

\section{Complete opcode contract}

After the four machines the architectural rectangle is fixed. The table is the single entry point for “which opcodes exist on which core”.

\begin{center}
\small
\begin{tabular}{@{}lcccc@{}}
\toprule
Class / feature & SVM-8 & SVM-16 & Zve32f & Zve64d/V \\
 & (BX-8) & (BX-16) & (BX-32) & (BX-64) \\
\midrule
\texttt{vsetvl}/\texttt{vsetvli} & yes & yes & yes & yes \\
\texttt{vsetivli} & opt. & yes & yes & yes \\
Unit-stride ld/st & EEW=8 & 8,16 & 8--32 & 8--64 \\
Constant-stride & --- & yes & yes & yes \\
Indexed ld/st & --- & --- & yes & yes \\
Segment ld/st & no & no & no & no \\
Integer arith (add/logic/shift) & yes & yes & yes & yes \\
min/max & --- & yes & yes & yes \\
Compare + mask logicals & yes & yes & yes & yes \\
Saturating add/sub & vsadd only & yes & yes & yes \\
Widening multiply / narrow & --- & yes & yes & yes \\
Reduction / slide1 & --- & yes & yes & yes \\
Gather / compress & --- & --- & yes & yes \\
Moves / splat (\texttt{vmv.*}) & yes & yes & yes & yes \\
FP32 vector + FMA & --- & --- & yes & yes \\
FP64 vector + FMA & --- & --- & --- & yes \\
Division / sqrt / estimates & no & no & opt. & per claim \\
LMUL $>1$ & no & widen only & yes & yes \\
Xsvm (custom-0) & no & opt. & opt. & off on V \\
Tail policy & und. & und. & und. & und. \\
\bottomrule
\end{tabular}
\end{center}

The bit-field diagrams that appear in Chapters~\ref{ch:bx8}--\ref{ch:bx64} are the concrete encodings for each class. The tables below list every mnemonic the corresponding profile must implement. The Bluespec package that follows elaborates exactly the cells marked “yes” (or “opt.” when the corresponding parameter is true) and turns every other cell into an illegal-instruction path.

\section{Aggregated opcode tables (ISA)}

Canonical bit layouts are shown once, width-aligned. Profile columns then list only the mnemonics (and the few constants that differ). Detailed per-machine diagrams remain in Chapters~\ref{ch:bx8}--\ref{ch:bx64}; this section is the compact academic summary.

\subsection{Canonical layouts (32-bit, two equal-width half-words)}

\paragraph{OP-V arithmetic / compare / move / FP}
\begin{center}
\ttfamily\footnotesize
\setlength{\tabcolsep}{0.3em}
\begin{tabular}{*{16}{w{c}{1.05em}}}
31 & 30 & 29 & 28 & 27 & 26 & 25 & 24 & 23 & 22 & 21 & 20 & 19 & 18 & 17 & 16 \\
\hline
\multicolumn{6}{|c|}{F} &
\multicolumn{1}{c|}{M} &
\multicolumn{5}{c|}{S} &
\multicolumn{4}{c|}{T} \\
\hline
\end{tabular}
\vspace{0.35ex}
\begin{tabular}{*{16}{w{c}{1.05em}}}
15 & 14 & 13 & 12 & 11 & 10 & 9 & 8 & 7 & 6 & 5 & 4 & 3 & 2 & 1 & 0 \\
\hline
\multicolumn{1}{|c|}{T} &
\multicolumn{3}{c|}{C} &
\multicolumn{5}{c|}{D} &
\multicolumn{7}{c|}{O} \\
\hline
\end{tabular}
\end{center}
\vspace{1.0ex}
\begin{center}\small
\begin{tabular}{@{}ll@{}}
F = funct6 & M = vm \\
S = vs2 & T = vs1 / rs1 / fs1 \\
C = funct3 & D = vd \\
O = 1010111 (OP-V) & \\
\end{tabular}
\end{center}

\paragraph{Unit-stride / stride / indexed load (store substitutes opcode 0100111)}
\begin{center}
\ttfamily\footnotesize
\setlength{\tabcolsep}{0.3em}
\begin{tabular}{*{16}{w{c}{1.05em}}}
31 & 30 & 29 & 28 & 27 & 26 & 25 & 24 & 23 & 22 & 21 & 20 & 19 & 18 & 17 & 16 \\
\hline
\multicolumn{3}{|c|}{N} &
\multicolumn{1}{c|}{E} &
\multicolumn{2}{c|}{P} &
\multicolumn{1}{c|}{M} &
\multicolumn{5}{c|}{L/S} &
\multicolumn{4}{c|}{B} \\
\hline
\end{tabular}
\vspace{0.35ex}
\begin{tabular}{*{16}{w{c}{1.05em}}}
15 & 14 & 13 & 12 & 11 & 10 & 9 & 8 & 7 & 6 & 5 & 4 & 3 & 2 & 1 & 0 \\
\hline
\multicolumn{3}{|c|}{W} &
\multicolumn{5}{c|}{D} &
\multicolumn{7}{c|}{O} \\
\hline
\end{tabular}
\end{center}
\vspace{1.0ex}
\begin{center}\small
\begin{tabular}{@{}ll@{}}
N = nf & E = mew \\
P = mop (00 unit, 10 stride, 01/11 indexed) & M = vm \\
L/S = lumop or vs2 (index) & B = rs1 (base) \\
W = width (000=8, 101=16, 110=32, 111=64) & D = vd / vs3 \\
O = 0000111 (load) & \\
\end{tabular}
\end{center}

\paragraph{Configuration (\texttt{vsetvli})}
\begin{center}
\ttfamily\footnotesize
\setlength{\tabcolsep}{0.3em}
\begin{tabular}{*{16}{w{c}{1.05em}}}
31 & 30 & 29 & 28 & 27 & 26 & 25 & 24 & 23 & 22 & 21 & 20 & 19 & 18 & 17 & 16 \\
\hline
\multicolumn{11}{|c|}{Z} &
\multicolumn{5}{c|}{S} \\
\hline
\end{tabular}
\vspace{0.35ex}
\begin{tabular}{*{16}{w{c}{1.05em}}}
15 & 14 & 13 & 12 & 11 & 10 & 9 & 8 & 7 & 6 & 5 & 4 & 3 & 2 & 1 & 0 \\
\hline
\multicolumn{3}{|c|}{C} &
\multicolumn{5}{c|}{D} &
\multicolumn{7}{c|}{O} \\
\hline
\end{tabular}
\end{center}
\vspace{1.0ex}
\begin{center}\small
\begin{tabular}{@{}ll@{}}
Z = zimm[10:0] & S = rs1 \\
C = 111 & D = rd \\
O = 1010111 & \\
\end{tabular}
\end{center}

\subsection{Mnemonic coverage by profile}

\begin{center}\small
\begin{tabular}{@{}lcccc@{}}
\toprule
Group & SVM-8 & SVM-16 & Zve32f & Zve64d/V \\
\midrule
Configuration & \texttt{vsetvl/i} & + \texttt{vsetivli} & same & same \\
Unit-stride & e8 & e8,e16 & e8--e32 & e8--e64 \\
Stride & --- & e8,e16 & e8--e32 & e8--e64 \\
Indexed & --- & --- & e8--e32 & e8--e64 \\
Integer arith & add/sub/logic/shift & +min/max & same & +SEW64 \\
Saturating & vsadd & full set & full & full \\
Widening / narrow & --- & vwmul, vnclip\ldots & same & same \\
Compare + mask & full & full & full & full \\
Reduction / slide1 & --- & yes & yes & yes \\
Gather / compress & --- & --- & yes & yes \\
Moves / splat & \texttt{vmv.*} & same & +FP moves & same \\
FP32 + FMA & --- & --- & full & full \\
FP64 + FMA & --- & --- & --- & full \\
Xsvm & no & opt. & opt. & off on V \\
\bottomrule
\end{tabular}
\end{center}

\subsection{Explicit summary --- SVM-8 (BX-8)}

\begin{center}\small
\begin{tabular}{@{}ll@{}}
\toprule
Required & Notes \\
\midrule
\texttt{vsetvl}, \texttt{vsetvli} & \texttt{vsetivli} optional \\
\texttt{vle8.v}, \texttt{vse8.v} & unit-stride only \\
\texttt{vadd/sub.vv/vx}, \texttt{vrsub.vx} & \\
\texttt{vand/vor/vxor.vv/vx} & \\
\texttt{vsll/vsrl/vsra.vv/vx} & \\
\texttt{vsadd.vv/vx} & \texttt{vssub} optional \\
\texttt{vmseq/ne/ltu/lt/leu/le.vv/vx}, \texttt{vmsgtu/gt.vx} & \\
mask logicals, \texttt{vcpop.m}, \texttt{vfirst.m} & \\
\texttt{vmv.v.v/x/i}, \texttt{vmv.x.s}, \texttt{vmv.s.x} & \\
\bottomrule
\end{tabular}
\end{center}

\subsection{Explicit summary --- SVM-16 (BX-16)}

Everything in SVM-8, plus:

\begin{center}\small
\begin{tabular}{@{}ll@{}}
\toprule
Required & Notes \\
\midrule
\texttt{vsetivli} & now required \\
\texttt{vle16.v}, \texttt{vse16.v} & \\
\texttt{vlse8/16.v}, \texttt{vsse8/16.v} & constant-stride \\
\texttt{vminu/min/maxu/max.vv/vx} & \\
\texttt{vsaddu/vsadd/vssubu/vssub} & full saturating set \\
\texttt{vwmul/u/su.vv/vx} & widening \\
\texttt{vnsrl/vnsra}, \texttt{vnclipu/vnclip} & narrowing \\
\texttt{vredsum.vs} & reduction \\
\texttt{vslide1up.vx}, \texttt{vslide1down.vx} & \\
Xsvm matrix / cross-product & optional (custom-0) \\
\bottomrule
\end{tabular}
\end{center}

\subsection{Explicit summary --- Zve32f (BX-32)}

Everything in SVM-16, plus:

\begin{center}\small
\begin{tabular}{@{}ll@{}}
\toprule
Required & Notes \\
\midrule
\texttt{vle32.v}, \texttt{vse32.v} & \\
\texttt{vluxei*}, \texttt{vloxei*}, indexed stores & \\
\texttt{vrgather.vv/vx}, \texttt{vcompress.vm} & \\
\texttt{vfadd/sub/mul.vv/vf} & FP32 \\
\texttt{vfmacc/fnmacc/fmsac/fnmsac} & FMA family \\
\texttt{vfmadd/fnmadd/fmsub/fnmsub} & \\
\texttt{vfmin/max}, \texttt{vfsqrt}, \texttt{vfsgnj*} & \\
\texttt{vfcvt.*}, \texttt{vmfeq/ne/lt/le/gt/ge} & \\
\texttt{vfmv.v.f}, \texttt{vfmv.f.s}, \texttt{vfmv.s.f} & \\
\bottomrule
\end{tabular}
\end{center}

\subsection{Explicit summary --- Zve64d / V (BX-64)}

Everything in Zve32f, plus SEW\,=\,64 and FP64. If the image claims full V, the entire V instruction set is required and the V bit is set in \texttt{misa}; otherwise the claim is Zve64d only. Xsvm must be off on any V-claiming image.

\subsection{Key constants (funct6 / width / mop)}

\begin{center}\small
\begin{tabular}{@{}lll@{}}
\toprule
Operation & Constant & Value \\
\midrule
\texttt{vadd} & funct6 & 000000 \\
\texttt{vsub} & funct6 & 000010 \\
\texttt{vand/vor/vxor} & funct6 & 001001/10/11 \\
\texttt{vsll/vsrl/vsra} & funct6 & 100101/101000/101001 \\
\texttt{vsadd} & funct6 & 100001 \\
\texttt{vwmul} & funct6 & 111011 \\
\texttt{vfmacc} & funct6 & 101100 \\
width 8/16/32/64 & W field & 000/101/110/111 \\
mop unit/stride/idx & P field & 00/10/01 or 11 \\
\bottomrule
\end{tabular}
\end{center}

Any mnemonic required by a wider profile but absent from a narrower one is an illegal-instruction trap on the narrower core. The bit layouts above are identical on every machine; only the legal SEW, the legal mop values, and the legal funct6 set change.

\section{Parameters}

The vector unit is one Bluespec package. The profile is an elaboration parameter, not a source fork.

\begin{lstlisting}
typedef struct {
  Integer elen;
  Integer vlen;
  Integer lanes;
  Bool    fp32;
  Bool    fp64;
  Bool    indexed;
  Bool    xsvm;
  Bool    claimV;
} VecParams deriving (Bits, Eq);
\end{lstlisting}

Legal combinations are exactly the four rows of the profile card. A function \texttt{paramsOk} returns false for every other combination, and the top-level module refuses to elaborate if it returns false. In particular \texttt{fp32} is false when \texttt{elen < 32}, and \texttt{claimV} is false when \texttt{elen < 64} or \texttt{vlen < 128} or \texttt{fp64} is false.

\section{Interfaces}

\begin{lstlisting}
interface VectorIssue;
  method Action req(Parcel p, Regs rs);
  method Bool ready;
endinterface

interface Lane;
  method Action enq(Elem a, Elem b, Opcode op, Bool mask);
  method ActionValue#(Elem) deq;
endinterface
\end{lstlisting}

The issue method is called from the hart's execute stage for OP-V parcels. \texttt{ready} is false while an in-order vector instruction occupies the unit. Loads use a third interface, \texttt{VecMem}, whose \texttt{beat} method is granted by the same arbiter as the video master on BX-8 and BX-16.

Scalar writeback of \texttt{vsetvl}'s destination register is a method on the hart, not a side path. The vector unit returns the granted \texttt{vl} in the result structure of \texttt{req}.

\section{The configuration rule}

One rule handles \texttt{vsetvl}, \texttt{vsetvli}, and, where the profile includes it, \texttt{vsetivli}. It computes the requested SEW and LMUL, compares them with \texttt{paramsOk} for this instruction, and either updates \texttt{vtype} and \texttt{vl} or sets \texttt{vill}. The rule is atomic, so a configuration cannot be observed half-written by a following vector instruction in the same schedule. Give the rule a higher urgency than the arithmetic rule so that a \texttt{vsetvl} decoded in cycle $n$ is visible to an instruction decoded in cycle $n+1$.

\section{Arithmetic rules}

Integer rules are always elaborated. Their bodies case on SEW and produce an illegal-instruction side effect when the SEW is outside the profile. That side effect is a second line of defence behind \texttt{vill}. The first line is the configuration rule. The second line catches an instruction that ignores \texttt{vtype}.

Floating-point rules are wrapped in \texttt{if (p.fp32)} blocks so that a BX-8 elaboration does not instantiate DSP blocks for a single-precision multiplier. The same pattern gates FP64, indexed address generation, and the Xsvm decoder.

\section{What SVM-8 still contains}

The SVM-8 elaboration still contains the decode for opcodes it will trap: the illegal path is a rule, and the rule is source text. That is intentional. A later profile turns the same decode into a data path by changing parameters, not by pasting a second decoder. Keep the trap path in the source even when you are sure the profile will never enable the opcode. The tests of Chapter~\ref{ch:tests} include an illegal SEW on every profile, and they need something in the design that raises it.
